% =========================
% MidTerm check in report
% =========================
\thispagestyle{empty}
\begin{center}
    {\Large \uline{Final Project Check-in Report}}
\end{center}

\subsection*{Updates based on feedback from Project Proposal Evaluation} The paper now has a more focused scope, concentrating on how generative AI tools reshape developer problem-solving strategies across different stages of programming tasks. The introduction has been revised to reflect this focus, emphasizing the cognitive and behavioral shifts induced by AI assistance. The title has been updated to "From Exploration to Acceleration: How Generative AI Reshapes Development Practices" to better capture the essence of the paper's focus on the transformation of development practices through generative AI.

\subsection*{(Updated) Scope of the Project}
This project investigates how generative AI tools influence developers’ problem-solving strategies during programming tasks. The scope of the work is to conduct a structured literature survey that synthesizes existing research on developer–AI interaction, with a particular focus on how these tools affect different stages of problem-solving. The key idea behind the project is to move beyond viewing AI tools as simple productivity enhancers and instead analyze their role in shaping cognitive processes and workflows. To achieve this, the paper adopts a hybrid approach that first defines a taxonomy of AI-assisted programming interactions and then maps these interactions to stages of software development such as ideation, implementation, review, and documentation. Compared to the original proposal, the scope has been refined to emphasize a structured taxonomy and stage-based analysis rather than a broader comparison of tools and outcomes. This change was made to provide a clearer analytical framework and enable deeper synthesis of insights across a limited set of papers.

\subsection*{Progress}
Significant progress has been made toward completing the project. A curated set of ten relevant papers has been identified and reviewed, covering empirical studies and surveys on generative AI in programming. The core structure of the paper, including the taxonomy of interaction modes and their mapping to development stages, has been established. Initial drafts of the introduction, motivation, taxonomy, and cross-stage analysis sections have been completed.

\subsection*{Challenges}
\begin{itemize}
    \item One challenge encountered was narrowing down the scope and selecting a manageable yet representative set of papers, given the rapidly growing body of literature in this area.
    \item Another challenge was synthesizing findings across studies with varying methodologies, participant populations, and AI tools, which required careful abstraction to identify common themes and patterns.
    \item (Optional) A third, yet optional challenge is procuring datasets that could help me replicate some of the findings from the papers, which would add an empirical dimension to the paper. \textit{This project is a literature survey at its core, but I am curious about the data that the original authors have used and how the findings relate to the broader context of the paper}, but is not essential for the core contributions.
\end{itemize}
Based on these challenges, the project timeline has been adjusted to allow for additional time in the synthesis and writing phases, particularly to ensure that the stage-based analysis is comprehensive and well-integrated with the taxonomy. The focus will be on deepening the analysis of how different interaction modes influence developer behavior at each stage of programming tasks, while also ensuring that the paper maintains a coherent narrative throughout.

\subsection*{Adjustments to the Project Plan}
Minor adjustments have been made to allocate additional time for synthesis and writing, particularly to ensure coherence between the taxonomy and stage-based analysis. These adjustments are manageable and do not affect the overall timeline for project completion.

\subsection*{Planned Work}
Based on the current progress, the next steps include:
\begin{itemize}
    \item Based on the depth of the topic, I have decided to include more research papers for review, trying to get an idea of what the landscape looked like before the recent surge in Generative AI and compare the benefits and drawbacks of Gen AI.
    \item I have \uline{excluded the discussion section for now}, which will be included in the final version after the stage-based analysis is complete. The discussion will synthesize the findings from the taxonomy and stage-based analysis, highlighting cross-cutting themes and implications for future research.
\end{itemize}

\subsection*{\textit{Updated} Timeline}
\begin{description}[leftmargin=!]
    \item[\checkmark \;Week 1:] Identified and finalized 10 papers
    \item[\checkmark \;Week 2:] Completed reading and initial notes
    \item[\checkmark \;Week 3:] Drafted taxonomy and interaction mapping
    \item[Week 4:] Writing stage-wise analysis and final sections
    \item[Week 5:] Collating all findings and create the Discussion section
    \item[Week 6:] Final revisions and submission
\end{description}

\noindent\rule{\columnwidth}{0.4pt}

\vspace{0.5em}

\textit{I have included the paper below as of the check-in report submission. This content is still a work in progress and may be revised in the \uline{final version}.}

\newpage